\section{P04 - Corrigé professeur - Types construits}

\subsection{Objectifs spécifiques}

\begin{itemize}
\tightlist
\item
  Objectif O1 - Identifier précisément la représentation ou la structure
  en jeu.
\item
  Objectif O2 - Appliquer une méthode disciplinaire complète.
\item
  Objectif O3 - Justifier le résultat sur un cas différent.
\item
  Objectif O4 - Contrôler un cas limite et corriger une erreur observée.
\end{itemize}

\subsection{Capacités officielles atomiques}

\begin{itemize}
\tightlist
\item
  P-DATA-CONSTR-02A
\item
  P-DATA-CONSTR-02D
\end{itemize}

\subsection{Prérequis}

\begin{itemize}
\tightlist
\item
  Reconnaître une consigne liée à tuple.
\item
  Distinguer donnée, méthode et conclusion dans le thème Tuples, listes,
  dictionnaires.
\item
  Rédiger une justification courte en utilisant le vocabulaire du
  programme.
\item
  Contrôler une réponse par un cas limite ou un contre-exemple
  explicite.
\end{itemize}

\subsection{Séance(s) correspondante(s)}

\begin{itemize}
\tightlist
\item
  P04-S1 à P04-S7 : support rattaché aux séances prêtes de la
  progression.
\end{itemize}

\subsection{Situation-problème concrète}

Une station météo stocke des coordonnées fixes, des relevés horaires
modifiables et des mesures accessibles par nom.

\subsection{Activité d’entrée}

\begin{enumerate}
\def\labelenumi{\arabic{enumi}.}
\tightlist
\item
  Identifier ce qui doit rester immuable dans un tuple.
\item
  Modifier une liste de températures.
\item
  Lire une clé dans un dictionnaire de station.
\item
  Décrire ce qui se passe avec une liste vide.
\end{enumerate}

\subsection{Méthode générale de correction}

\begin{itemize}
\tightlist
\item
  Point 1 : pour tuple de coordonnées, exiger la donnée
  \texttt{(36.8,\ 10.2)}, la méthode « lire sans modifier et nommer
  latitude puis longitude » et le contrôle « tentative de modification
  interdite ».
\item
  Point 2 : pour liste de relevés, exiger la donnée
  \texttt{{[}18,\ 20,\ 19{]}}, la méthode « parcourir les valeurs et
  calculer une moyenne » et le contrôle « liste vide ».
\item
  Point 3 : pour dictionnaire, exiger la donnée
  \texttt{\{"temperature":\ 21,\ "vent":\ 12\}}, la méthode « tester la
  présence de la clé avant lecture » et le contrôle « clé absente ».
\item
  Point 4 : pour copie de liste, exiger la donnée
  \texttt{{[}{[}1{]},\ {[}2{]}{]}}, la méthode « distinguer copie
  superficielle et copie indépendante » et le contrôle « liste imbriquée
  ». \#\# Exercices numérotés
\item
  Exercice 1 : résoudre tuple de coordonnées avec \texttt{(36.8,\ 10.2)}
  ; attendu : coordonnées conservées.
\item
  Exercice 2 : expliquer liste de relevés à partir de
  \texttt{{[}18,\ 20,\ 19{]}} ; attendu : \texttt{19}.
\item
  Exercice 3 : comparer dictionnaire avec
  \texttt{\{"temperature":\ 21,\ "vent":\ 12\}} ; attendu : \texttt{21}
  pour \texttt{temperature}.
\item
  Exercice 4 : corriger copie de liste pour
  \texttt{{[}{[}1{]},\ {[}2{]}{]}} ; attendu : modification locale
  contrôlée.
\item
  Exercice 5 : tester un cas limite lié à tentative de modification
  interdite ; attendu : le comportement de tuple de coordonnées est
  contrôlé.
\item
  Exercice 6 : classer deux méthodes possibles pour liste de relevés ;
  attendu : la méthode robuste est choisie et justifiée.
\item
  Exercice 7 : justifier un transfert qui utilise dictionnaire avec une
  donnée nouvelle ; attendu : la justification reste valable sur le
  nouveau cas.
\item
  Exercice 8 : étendre un énoncé volontairement erroné sur copie de
  liste ; attendu : l'erreur est localisée puis réparée.
\end{itemize}

\subsection{Corrigé}

\subsubsection{Corrigé exercice 1}

\begin{itemize}
\tightlist
\item
  Méthode : identifier \texttt{(36.8,\ 10.2)}, appliquer la méthode «
  lire sans modifier et nommer latitude puis longitude », puis écrire
  coordonnées conservées.
\item
  Résultat : coordonnées conservées.
\item
  Contrôle : faire apparaître le contrôle « tentative de modification
  interdite ».
\item
  Erreur traitée : EF1 - Modifier un tuple comme une liste. \#\#\#
  Corrigé exercice 2
\item
  Méthode : expliciter chaque étape de parcourir les valeurs et calculer
  une moyenne avant de conclure par \texttt{19}.
\item
  Résultat : \texttt{19}.
\item
  Contrôle : rédiger la méthode avant le résultat.
\item
  Erreur traitée : EF2 - Parcourir les indices quand les valeurs
  suffisent. \#\#\# Corrigé exercice 3
\item
  Méthode : comparer la donnée avec le cas limite « clé absente » et
  valider \texttt{21} pour \texttt{temperature}.
\item
  Résultat : \texttt{21} pour \texttt{temperature}.
\item
  Contrôle : comparer avec le cas « clé absente ».
\item
  Erreur traitée : EF3 - Accéder à une clé sans vérifier sa présence.
  \#\#\# Corrigé exercice 4
\item
  Méthode : isoler l'erreur fréquente « Copier une liste imbriquée
  seulement au premier niveau. » puis reprendre la procédure correcte.
\item
  Résultat : modification locale contrôlée.
\item
  Contrôle : corriger l'erreur « Copier une liste imbriquée seulement au
  premier niveau. ».
\item
  Erreur traitée : EF4 - Copier une liste imbriquée seulement au premier
  niveau. \#\#\# Corrigé exercice 5
\item
  Méthode : identifier \texttt{(36.8,\ 10.2)}, appliquer la méthode «
  lire sans modifier et nommer latitude puis longitude », puis écrire
  coordonnées conservées.
\item
  Résultat : le comportement de tuple de coordonnées est contrôlé.
\item
  Contrôle : nommer la donnée minimale et la conclusion.
\item
  Erreur traitée : EF1 - Modifier un tuple comme une liste. \#\#\#
  Corrigé exercice 6
\item
  Méthode : expliciter chaque étape de parcourir les valeurs et calculer
  une moyenne avant de conclure par \texttt{19}.
\item
  Résultat : la méthode robuste est choisie et justifiée.
\item
  Contrôle : identifier pourquoi « Parcourir les indices quand les
  valeurs suffisent. » est une erreur.
\item
  Erreur traitée : EF2 - Parcourir les indices quand les valeurs
  suffisent. \#\#\# Corrigé exercice 7
\item
  Méthode : comparer la donnée avec le cas limite « clé absente » et
  valider \texttt{21} pour \texttt{temperature}.
\item
  Résultat : la justification reste valable sur le nouveau cas.
\item
  Contrôle : inclure une étape calculable par un pair.
\item
  Erreur traitée : EF3 - Accéder à une clé sans vérifier sa présence.
  \#\#\# Corrigé exercice 8
\item
  Méthode : isoler l'erreur fréquente « Copier une liste imbriquée
  seulement au premier niveau. » puis reprendre la procédure correcte.
\item
  Résultat : l'erreur est localisée puis réparée.
\item
  Contrôle : proposer une activité corrective inspirée de « Modifier une
  sous-liste et observer l'effet sur la copie. ».
\item
  Erreur traitée : EF4 - Copier une liste imbriquée seulement au premier
  niveau.
\end{itemize}

\subsection{Barème de correction rapide}

\begin{itemize}
\tightlist
\item
  Exercice 1 : 1 point méthode, 0,5 point résultat, 0,5 point contrôle
  sur « faire apparaître le contrôle « tentative de modification
  interdite » ».
\item
  Exercice 2 : 1 point méthode, 0,5 point résultat, 0,5 point contrôle
  sur « rédiger la méthode avant le résultat ».
\item
  Exercice 3 : 1 point méthode, 0,5 point résultat, 0,5 point contrôle
  sur « comparer avec le cas « clé absente » ».
\item
  Exercice 4 : 1 point méthode, 0,5 point résultat, 0,5 point contrôle
  sur « corriger l'erreur « Copier une liste imbriquée seulement au
  premier niveau. » ».
\item
  Exercice 5 : 1 point méthode, 0,5 point résultat, 0,5 point contrôle
  sur « nommer la donnée minimale et la conclusion ».
\item
  Exercice 6 : 1 point méthode, 0,5 point résultat, 0,5 point contrôle
  sur « identifier pourquoi « Parcourir les indices quand les valeurs
  suffisent. » est une erreur ».
\item
  Exercice 7 : 1 point méthode, 0,5 point résultat, 0,5 point contrôle
  sur « inclure une étape calculable par un pair ».
\item
  Exercice 8 : 1 point méthode, 0,5 point résultat, 0,5 point contrôle
  sur « proposer une activité corrective inspirée de « Modifier une
  sous-liste et observer l'effet sur la copie. » ». \#\# Erreurs
  fréquentes
\item
  Erreur fréquente EF1 - Modifier un tuple comme une liste.
\item
  Erreur fréquente EF2 - Parcourir les indices quand les valeurs
  suffisent.
\item
  Erreur fréquente EF3 - Accéder à une clé sans vérifier sa présence.
\item
  Erreur fréquente EF4 - Copier une liste imbriquée seulement au premier
  niveau.
\end{itemize}

\subsection{Remédiation ciblée}

\begin{itemize}
\tightlist
\item
  Activité corrective EF1 : Identifier mutabilité et usage avant
  d'écrire une affectation.
\item
  Activité corrective EF2 : Écrire deux boucles, avec indices puis avec
  valeurs, et comparer.
\item
  Activité corrective EF3 : Tester \texttt{cle\ in\ dictionnaire} avant
  la lecture.
\item
  Activité corrective EF4 : Modifier une sous-liste et observer l'effet
  sur la copie.
\end{itemize}

\subsection{Différenciation}

\begin{itemize}
\tightlist
\item
  Socle : traiter \texttt{(36.8,\ 10.2)} avec une fiche méthode fournie.
\item
  Standard : traiter \texttt{{[}18,\ 20,\ 19{]}} en rédigeant la
  justification complète.
\item
  Expert : inventer un cas limite lié à « clé absente » et expliquer le
  comportement attendu.
\end{itemize}

\subsection{Critères de réussite}

\begin{itemize}
\tightlist
\item
  La capacité officielle est citée dans la copie.
\item
  La méthode contient au moins une étape vérifiable par un pair.
\item
  Le cas limite est discuté avec une donnée concrète.
\item
  La correction explique quelle erreur fréquente est évitée.
\end{itemize}

\subsection{Corrigés complémentaires - types construits}

\subsubsection{Corrigé exercice 9}

\begin{itemize}
\tightlist
\item
  Valeur attendue : \texttt{A=(0,4)}, \texttt{B=(6,10)}, milieu
  \texttt{(3.0,\ 7.0)}.
\item
  Justification : la formule moyenne les abscisses et les ordonnées
  séparément.
\item
  Contrôle : le résultat est encore un tuple de taille 2.
\end{itemize}

\subsubsection{Corrigé exercice 10}

\begin{itemize}
\tightlist
\item
  Valeur attendue : \texttt{{[}8,\ 14,\ 10{]}}.
\item
  Justification : une liste est mutable ; l'affectation
  \texttt{notes{[}1{]}\ =\ 14} modifie la liste en place.
\item
  Contrôle : l'indice \texttt{1} désigne la deuxième valeur, pas la
  première.
\end{itemize}

\subsubsection{Corrigé exercice 11}

\begin{itemize}
\tightlist
\item
  Valeur attendue : \texttt{{[}"A",\ "C"{]}}.
\item
  Justification : les dictionnaires des stations A et C ont
  \texttt{temperature\ \textgreater{}=\ 20}.
\item
  Contrôle : le filtre lit la clé \texttt{temperature} et ne compare pas
  le dictionnaire complet.
\end{itemize}

\subsubsection{Corrigé exercice 12}

\begin{itemize}
\tightlist
\item
  Valeurs attendues : mauvaise taille pour \texttt{(2,)}, clé absente
  pour \texttt{"pression"}, liste vide pour \texttt{{[}{]}}.
\item
  Justification : chaque structure impose une vérification différente
  avant traitement.
\item
  Contrôle : le corrigé cite l'erreur et la garde à ajouter.
\end{itemize}
